\labsection{4}{CPU scheduling algorithms}{sec:lab04}
\begin{labproject}{Trace and compare four CPU schedulers}
\textbf{Coverage.} Chapter 4.\\
\textbf{Source files.} Four C programs in \texttt{Lab04}.\\
\textbf{Goal.} Implement and compare FCFS, non-preemptive SJF, priority scheduling, and Round Robin.

Use one common workload so comparisons are meaningful. For the supplied programs, arrival times are effectively assumed equal unless you extend the code.

\begin{mltable}
\begin{tabularx}{\linewidth}{@{}>{\bfseries\leavevmode\color{MLNavy}}p{0.20\linewidth}p{0.23\linewidth}X@{}}
\toprule
\rowcolor{MLPaleBlue!92}
Program & Policy & Verify \\
\midrule
\texttt{Task03.c} & FCFS & cumulative waiting time in input order \\
\texttt{Task04.c} & SJF & ascending burst order and original PID preservation \\
\texttt{Task01.c} & Priority & priority convention and tie handling \\
\texttt{Task02.c} & Round Robin & quantum, remaining burst, completion, waiting calculation \\
\bottomrule
\end{tabularx}
\end{mltable}

\begin{debugtask}
For Round Robin, validate waiting time against a hand-drawn Gantt chart. A program can print plausible totals while still updating waiting time incorrectly during intermediate quanta.
\end{debugtask}

\begin{tryit}
Use bursts P1=8, P2=4, P3=2, P4=1. Compare FCFS and SJF. Then use RR with quantum 2 and report first-response time for every process.
\end{tryit}
\end{labproject}

\begin{labdeliverable}
Submit source/output for all four algorithms, one common workload, hand-checked Gantt charts, and a comparison table of average waiting and turnaround times.
\end{labdeliverable}
